\section{Worked examples}
\label{sec:examples_worked}

Throughout we compute in the configuration–space model with kernels built from the HT propagator $K=K^{\C}\otimes K^{\R}$:
$K^{\C}$ has a simple logarithmic pole along diagonals in $\FM_2(\C)$ and $K^{\R}$ is the standard step‑function (time‑ordering) kernel on~$\R$ (or the half‑line for boundary computations).

\subsection{Free chiral multiplet}
\label{subsec:free-chiral}
Consider a free chiral multiplet in the HT twist. The fields are $(\phi,\bar\phi)$ and their superpartners with BV completions. The action is quadratic; hence all $m_{k\ge3}$ vanish by Wick's theorem.

\paragraph{Binary operation.}
The two‑point OPE kernel is
\[
m_2(a,b)(z) \;=\; \int_{\FM_2(\C)} K^{\C}(z_{12})\cdot a\otimes b,
\]
with $K^{\C}(z_{12})\sim \frac{1}{z_{12}}$ in the holomorphic coordinate. The regular part (nonnegative powers of $\lambda$) yields a strictly commutative product on cohomology; the antisymmetric singular part vanishes in cohomology by exactness of $\dbar$ on the kernel.

\begin{proposition}[PVA for the free multiplet; \ClaimStatusConditional]
\label{prop:PVA-free}
Under Hypotheses~\textrm{(H1)--(H4)}, $H^\bullet(A,Q)$ is a commutative vertex algebra; the $(-1)$–shifted Poisson bracket on cohomology is trivial.
\end{proposition}

\begin{proof}
Immediate from the fact that $m_2^{\mathrm{sing}}$ is $Q$–exact (the residue cancels after integrating the distributional kernel against closed test fields) and $m_{k\ge3}=0$.
\end{proof}

\subsection{Landau–Ginzburg model with cubic superpotential}
\label{subsec:LG-cubic}
Consider a single chiral field $\phi$ with superpotential $W(\phi)=\frac{\lambda}{3}\phi^3$.
In the HT twist, the interaction vertex induces a nontrivial $m_3$ at tree level.

\paragraph{The $m_3$ integral.}
At leading order in $\lambda$,
\[
m_3(a,b,c) \;=\; \lambda \int_{\FM_3(\C)} K^{\C}(z_{12})K^{\C}(z_{23})K^{\C}(z_{31}) \cdot (a\otimes b\otimes c),
\]
with ordered times in $\Conf_3(\R)$ capturing the $E_1$ homotopy; the total kernel has degree $-2$ (cohomological).  Boundary faces of $\FM_3(\C)$ encode the Stasheff identity at arity $3$.

\paragraph{Cohomology.}
The $m_3$ appears only at chain level; using Stokes' theorem and AOS, its contribution to associativity is $Q$–exact, hence vanishes on $H^\bullet$. However, the antisymmetric part of $m_2$ produces a nontrivial $(-1)$–shifted bracket on $H^\bullet$ proportional to~$\lambda$.

\begin{proposition}[PVA for LG–cubic; \ClaimStatusConditional]
\label{prop:PVA-LG}
Under Hypotheses~\textrm{(H1)--(H4)}, $H^\bullet(A,Q)$ is a $(-1)$–shifted PVA with nontrivial bracket $\{\phi_\lambda \phi\}\sim \lambda\, \lambda^0$ (up to normalization), while $m_{k\ge3}$ vanish in cohomology.
\end{proposition}

\begin{proof}
The bracket originates from the singular residue of $m_2$; $m_3$ contributes homotopies only. Computation follows from the standard $\FM_3(\C)$ integral and the LG interaction vertex.
\end{proof}

\subsection{Abelian Chern–Simons with boundary}
\label{subsec:CS-abelian}
Let $A$ be abelian Chern–Simons theory on a half–space with HT boundary condition. The boundary algebra contains a current $J(z)$ with OPE
\[
J(z)J(w) \;\sim\; \frac{k}{(z-w)^2} + \text{regular},
\]
at level $k$ determined by the CS coupling.

\paragraph{Bulk-to-boundary and $R(z)$.}
Place two line operators of charges $q_1,q_2$ at $z_1,z_2$ on the boundary. The bulk–to–boundary kernel yields a spectral braiding
\[
R(z)\cdot (v_{q_1}\otimes v_{q_2}) \;=\; \exp\!\Bigl( \hbar\, \frac{q_1 q_2}{z} \Bigr) \,(v_{q_2}\otimes v_{q_1}),
\]
up to normalization (in perturbation theory). The classical $r(z)=\frac{q_1 q_2}{z}$ satisfies CYBE.

\begin{proposition}[Boundary Kac–Moody and spectral $R$; \ClaimStatusConditional]
\label{prop:CS-R}
Under Hypotheses~\textrm{(H1)--(H4)}, the boundary category of line operators is braided by $R(z)$, with classical limit $r(z)=\frac{k}{z}$ (after normalization). The bulk chiral Hochschild cochains produce the same $r(z)$ via the Laplace transform of the bulk $\lambda$–bracket kernel as in Proposition~\ref{prop:field-theory-r}.
\end{proposition}

\begin{proof}
The current algebra OPE computes the leading residue; the configuration integral over $\FM_2(\C)$ reproduces $1/z$. The Laplace transform calculation matches the boundary braiding.
\end{proof}

\subsection{DK-0 Laplace verification: the standard landscape}

\begin{proposition}[DK-0 Laplace bridge for five families; \ClaimStatusProvedHere]
\label{prop:dk0-laplace-five-families}
For each of the five standard families
$\cA \in \{\mathcal{H}_k,\;
\widehat{\mathfrak{sl}}_2,\;
\mathrm{Vir}_c,\;
\beta\gamma\text{/}bc,\;
\mathcal{W}_3\}$,
the Laplace transform of the $\lambda$-bracket equals the classical
$r$-matrix extracted from OPE singular terms:
\[
  r^{ab}(z)
  \;=\;
  \sum_{n \geq 0}
  c_n^{ab} \cdot \frac{n!}{z^{n+1}},
  \qquad
  \text{where }
  \{a_\lambda b\} = \sum_{n \geq 0} c_n^{ab}\, \lambda^n.
\]
Moreover, the modular characteristic $\kappa(\cA) + \kappa(\cA^!)$
is a constant independent of the level/charge parameter for each family,
and Koszul duality acts on the chiral Hochschild complex as predicted by
Theorem~H of Volume~I.
\end{proposition}

\begin{proof}
Direct computation on each family.  The Laplace identity
is Proposition~\ref{prop:field-theory-r} specialized to explicit
$\lambda$-bracket data.
For Heisenberg: $\{a_\lambda a\} = k\lambda$ gives
$r(z) = k/z^2$.
For $\widehat{\mathfrak{sl}}_2$: all $9$ bracket pairs
$(J^a, J^b)$ verify.
For Virasoro: $\{T_\lambda T\} = \partial T + 2T\lambda
+ (c/12)\lambda^3$ gives $r(z) = \partial T/z + 2T/z^2
+ (c/2)/z^4$, matching the OPE.
For $\beta\gamma$/$bc$: $\{b_\lambda c\} = 1$ gives $r(z) = 1/z$.
For $\mathcal{W}_3$: all three bracket pairs verify.
Complementarity constants: $0$ (Heisenberg, $\mathfrak{sl}_2$,
$bc$), $13$ (Virasoro), $250/3$ ($\mathcal{W}_3$).
Computational verification: 26 checks across 5 bridges in
\texttt{compute/lib/cross\_volume\_bridge.py}.
\end{proof}

\subsection{Kodaira--Spencer universal defect algebra}
\label{subsec:KS-defect}

In Zeng's boundary-algebra formulation of twisted holography~\cite{Zeng2024,CDG2023},
the large-$N$ open-side chiral algebra is built from the following field content
in $\mathfrak{gl}_N$: a $(b,c)$ ghost system in the adjoint,
a pair $(Z_1,Z_2)$ of adjoint $\beta\gamma$ systems, and
a pair $(I,J)$ of fundamental/anti-fundamental $\beta\gamma$ systems.

\paragraph{BRST charge.}
The BRST differential is
\[
  Q \;=\; \oint\!\bigl(\operatorname{Tr}(c\,Z_1 Z_2)
  + \operatorname{Tr}(I\,c\,J)
  + \tfrac{1}{2}\operatorname{Tr}(b\,c^2)\bigr).
\]
The first two terms encode the holomorphic moment map constraints
(complex Chern--Simons on $\C^3$ after the holomorphic--topological twist),
while the ghost self-coupling $\frac{1}{2}\operatorname{Tr}(bc^2)$
implements the BRST closure $Q^2=0$ via the Maurer--Cartan equation
in the $\mathfrak{gl}_N$ gauge sector.

\paragraph{Large-$N$ cohomology generators.}
At $N=\infty$ the BRST cohomology is freely generated by
single-trace operators
\[
  A(n),\; B(n),\; C(n),\; D(n),\; E_t(n),
  \qquad n \ge 0,
\]
where $A(n)=\operatorname{Tr}(Z_1^n\,\partial Z_1)$, etc., and $E_t(n)$
denotes the bilinear $I$--$J$ family labelled by a spectral parameter~$t$.
These generate a vertex algebra isomorphic to the Kodaira--Spencer
chiral algebra of~\cite{BCOV93}, extended by the defect sector.

\paragraph{The universal defect algebra.}
Define $U_{\mathrm{KS}} := A_{\partial,\mathrm{KS}}(B_{\mathrm{univ}})$, the
universal defect algebra obtained by applying the boundary
bar construction of Part~I to the universal bulk input~$B_{\mathrm{univ}}$.
The KK (Kaluza--Klein) reduction mechanism identifies
\[
  \text{boundary chiral algebra}
  \;=\;
  \text{universal defect chiral algebra of the parent 6d theory},
\]
providing a bridge from the large-$N$ gauged $\beta\gamma$ VOA
to Kodaira--Spencer gravity.  In bar-cobar language:
$U_{\mathrm{KS}} \simeq \Omega B(\cA_{\mathrm{open}})$ on the Koszul locus,
so bar-cobar inversion (Theorem~B of Volume~I) controls the
passage from open-string to closed-string (gravitational) degrees of freedom.

\subsection{The M2-brane instanton particle: bar-cobar as holographic duality}
\label{subsec:M2-holography}

Costello's M2-brane program~\cite{Costello-1705.02500v1,costello-gaiotto}
identifies, at large rank~$K$, the M2-brane boundary algebra
as Koszul dual to the 11d twisted supergravity bulk algebra.
This furnishes the most constructive realization of the thesis that
\emph{bar-cobar duality IS holographic duality}.

\paragraph{Seed algebra.}
The classical starting point is the double-current algebra
\[
  \mathfrak{g}_{\mathrm{dbl}}
  \;=\; \mathfrak{gl}_r \otimes \C[u,v],
\]
a polynomial-current extension of $\mathfrak{gl}_r$ in two spectral variables.
Its enveloping algebra $U(\mathfrak{g}_{\mathrm{dbl}})$ carries a natural
filtration by polynomial degree in $(u,v)$.

\paragraph{Koszul dual.}
The Koszul dual of the M2-brane algebra is computed by the bar-cobar construction:
\[
  A^!_{\mathrm{M2}}
  \;=\; \Omega B\bigl(U(\mathfrak{g}_{\mathrm{dbl}})\bigr)
  \;\simeq\; C^\bullet(\mathfrak{g}_{\mathrm{dbl}}),
\]
where $C^\bullet(\mathfrak{g}_{\mathrm{dbl}})$ denotes the Chevalley--Eilenberg
cochains.  The quasi-isomorphism follows from Theorem~B (bar-cobar inversion)
applied on the Koszul locus, which is verified for $U(\mathfrak{g}_{\mathrm{dbl}})$
by the PBW filtration and the classical Koszulness of $\operatorname{Sym}$.

\paragraph{Sector-by-sector matching.}
For each sector~$\sigma$ and level~$\ell$, the holographic dictionary reads
\[
  \bigl(A_{\partial,\infty}(\sigma,\ell)\bigr)^!
  \;\cong\;
  A_{\mathrm{bulk}}(\sigma,\ell),
\]
where $A_{\partial,\infty}$ is the large-$K$ boundary algebra and
$A_{\mathrm{bulk}}$ is the twisted supergravity bulk algebra.
The Koszul dual functor $(-)^!$ exchanges the boundary's associative
product with the bulk's Lie bracket, exactly as $\mathrm{Com}^!=\mathrm{Lie}$
at the operadic level.

\paragraph{Line operators and spectral kernel.}
The category of M2-brane line operators admits a description
\[
  \cC_{\mathrm{M2}}
  \;\simeq\;
  A^!_{\mathrm{bulk}}\text{-}\mathrm{mod},
  \qquad
  \text{with spectral kernel } r_{\mathrm{M2}}(u),
\]
where $r_{\mathrm{M2}}(u)$ is the classical $r$-matrix obtained from the
Laplace transform of the $\lambda$-bracket on $\mathfrak{g}_{\mathrm{dbl}}$,
as in Proposition~\ref{prop:field-theory-r}.  The spectral braiding
reproduces the RTT relations of the shifted Yangian~\cite{DNP2025}.
This closes the circle: bar-cobar duality at the algebraic level
\emph{is} holographic duality at the physical level.


